\chapter{Chaînes de Markov et files d'attente}\label{ch:qm:markov-chains-and-queues}

Au meilleur cours acheteur d'un contrat à terme à large tick, 800 lots sont en attente~; au meilleur cours vendeur, 200. Un teneur de marché vient de
s'insérer en queue de la file acheteuse et veut connaître deux choses~: la probabilité que son ordre soit exécuté
avant que la file vendeuse ne soit consommée et que le prix ne monte d'un tick, laissant l'ordre sans exécution~; et combien de temps
il devra attendre. Les deux questions portent sur le temps qu'une marche aléatoire met pour atteindre zéro. Si les ordres entrent et
sortent d'une file à des instants aléatoires, chaque file est une \omterm{def:qm:markov-chains-and-queues:ctmc}{chaîne de Markov}, et les réponses se déduisent de
la \omterm{def:qm:markov-chains-and-queues:generator}{matrice génératrice} de la chaîne~: 73~\% et 82 secondes, pour les taux de ce chapitre, contre 99~\% si la file vendeuse
était deux fois plus longue. Ce chapitre présente les \omterm{def:qm:markov-chains-and-queues:ctmc}{chaînes de Markov à temps continu}, les \omterm{def:qm:markov-chains-and-queues:bd}{processus de naissance et de mort}
et les files simples qui en découlent, l'\omterm{met:qm:markov-chains-and-queues:firststep}{analyse au premier pas} qui calcule les \omterm{def:qm:markov-chains-and-queues:hitting}{probabilités d'atteinte} et
les \omterm{def:qm:brownian-motion:passage}{temps d'atteinte}, ainsi qu'un modèle des deux meilleures files du carnet d'ordres, dans l'esprit de Cont et de Larrard
(2013).

\section{Chaînes de Markov à temps continu}

\begin{definition}[Chaîne de Markov, chaîne de Markov à temps continu]\label{def:qm:markov-chains-and-queues:ctmc}
Une \emph{chaîne de Markov}\index{chaîne de Markov} est un \omterm{def:qm:brownian-motion:markov}{processus de Markov} à espace d'états dénombrable, observé
à temps entiers, de matrice de transition $P_{ij} = \P(X_{n+1} = j \mid X_n = i)$. Une
\emph{chaîne de Markov à temps continu}\index{chaîne de Markov à temps continu} est un \omterm{def:qm:brownian-motion:markov}{processus de Markov} continu à droite
$(X_t)_{t\ge0}$ à espace d'états dénombrable et dont les probabilités de transition $P_{ij}(t) = \P(X_{s+t}
= j \mid X_s = i)$ ne dépendent pas de $s$.
\end{definition}

\begin{definition}[Matrice génératrice, chaîne des sauts]\label{def:qm:markov-chains-and-queues:generator}
La \emph{matrice génératrice}\index{matrice génératrice} d'une chaîne à temps continu est $\mathcal L = (q_{ij})$ avec $q_{ij}
= \lim_{t\downarrow0}P_{ij}(t)/t \ge 0$ pour $j \ne i$, le taux de sauts de $i$ vers $j$, et $q_{ii} = -\sum_{j\ne i}q_{ij}$. La
\emph{chaîne des sauts}\index{chaîne des sauts} est la chaîne à temps discret des états visités, dont les probabilités de transition
$q_{ij}/q_i$, où $q_i = -q_{ii}$.
\end{definition}

\begin{proposition}[Construction et équations de Kolmogorov]\label{prop:qm:markov-chains-and-queues:kolmogorov}
Une chaîne finie de \omterm{def:qm:markov-chains-and-queues:generator}{matrice génératrice} $\mathcal L$ demeure dans l'état $i$ pendant une durée exponentielle de taux $q_i$, puis se déplace
selon la \omterm{def:qm:markov-chains-and-queues:generator}{chaîne des sauts}, indépendamment du passé. Sa matrice de transition est $P(t) = e^{t\mathcal L}$, qui
vérifie $P'(t) = P(t)\mathcal L = \mathcal LP(t)$, les équations progressive et rétrograde.
\end{proposition}

\begin{proof}
La durée de séjour exponentielle est la seule loi sans mémoire, que la propriété de Markov impose à la durée
passée dans un état. Pour $h$ petit, $P(t + h) = P(t)P(h) = P(t)(I + h\mathcal L + o(h))$~; l'exponentielle matricielle est la
solution unique telle que $P(0) = I$.
\end{proof}

Une chaîne à temps continu est, pour un espace d'états fini, le pendant du générateur discret du chapitre~4~: les
mêmes équations progressive et rétrograde, avec une matrice à la place d'un opérateur différentiel. Sa \omterm{def:qm:stochastic-differential-equations:kolmogorov}{loi
stationnaire} vérifie $\pi\mathcal L = 0$.

\begin{example}[La fourchette en ticks]\label{ex:qm:markov-chains-and-queues:spread}
Modélisons la fourchette d'un contrat à large tick comme valant 1, 2 ou 3 ticks, qui s'élargit de 1 à 2 au taux de 2 par seconde et de 2
à 3 au taux de 1, se resserre de 2 à 1 au taux de 5, de 3 à 2 au taux de 4 et de 3 à 1 au taux de 0.5. En résolvant $\pi\mathcal L
= 0$ on obtient $\pi = (0.676, 0.265, 0.059)$, soit une fourchette moyenne de 1.38 ticks~; la fourchette reste à un tick pendant 0.5 seconde
en moyenne avant de s'élargir.
\end{example}

\begin{definition}[Équilibre détaillé]\label{def:qm:markov-chains-and-queues:balance}
Une loi $\pi$ vérifie l'\emph{équilibre détaillé}\index{equilibre@équilibre détaillé} pour un générateur donné si $\pi_iq_{ij} =
\pi_jq_{ji}$ pour tous $i, j$~: en régime stationnaire, chaque transition est aussi fréquente que sa réciproque, et la chaîne parcourue
à rebours dans le temps a la même loi (elle est réversible).
\end{definition}

L'\omterm{def:qm:markov-chains-and-queues:balance}{équilibre détaillé} implique $\pi\mathcal L = 0$ en sommant sur $i$~; la réciproque est fausse lorsque la probabilité circule
autour d'un cycle, comme dans l'exemple de la fourchette (l'\cref{exo:qm:markov-chains-and-queues:4}). Les chaînes de naissance et de mort,
qui ne se déplacent qu'un pas à la fois, vérifient toujours cette propriété.

\section{Processus de naissance et de mort et files simples}

\begin{definition}[Processus de naissance et de mort]\label{def:qm:markov-chains-and-queues:bd}
Un \emph{processus de naissance et de mort}\index{processus de naissance et de mort} est une chaîne à temps continu sur $\{0, 1, 2, \dots\}$ dont les seuls
sauts sont $n \to n + 1$ au taux $\lambda_n$ et $n \to n - 1$ au taux $\mu_n$.
\end{definition}

\begin{proposition}[Loi stationnaire d'un processus de naissance et de mort]\label{prop:qm:markov-chains-and-queues:bdstat}
Si $\sum_n\prod_{k=1}^n\lambda_{k-1}/\mu_k < \infty$, la \omterm{def:qm:stochastic-differential-equations:kolmogorov}{loi stationnaire} est $\pi_n \propto \prod_{k=1}^n\lambda_{k-1}/\mu_k$.
\end{proposition}

\begin{proof}
L'\omterm{def:qm:markov-chains-and-queues:balance}{équilibre détaillé} entre $n - 1$ et $n$ s'écrit $\pi_{n-1}\lambda_{n-1} = \pi_n\mu_n$~; on itère puis on normalise.
\end{proof}

\begin{definition}[File M/M/1]\label{def:qm:markov-chains-and-queues:mm1}
La \emph{file M/M/1}\index{file M/M/1} a des arrivées poissonniennes au taux $\lambda$, des temps de service exponentiels de
taux $\mu$ et un seul serveur~; le nombre d'éléments présents dans le système est un \omterm{def:qm:markov-chains-and-queues:bd}{processus de naissance et de mort} avec $\lambda_n = \lambda$ et $\mu_n =
\mu$. (notation de Kendall~: arrivées markoviennes, service markovien, un seul serveur.)
\end{definition}

Avec un taux d'utilisation $\rho = \lambda/\mu < 1$, la \omterm{def:qm:stochastic-differential-equations:kolmogorov}{loi stationnaire} est géométrique, $\pi_n = (1 - \rho)\rho^n$, si bien que le nombre moyen dans le
système est $L = \rho/(1 - \rho)$. La loi de Little, $L = \lambda W$ pour le temps moyen $W$ passé dans le système, vaut
pratiquement toute file stationnaire (Little, 1961), et donne $W = 1/(\mu - \lambda)$. Les deux explosent lorsque $\rho \to 1$~: un moteur d'appariement
ou une passerelle de contrôle des risques qui fonctionne à 90~\% de sa capacité compte en moyenne neuf messages en attente, à 95~\%
dix-neuf (\cref{fig:qm:markov-chains-and-queues:mm1}). Le livre~13 dimensionne ses files avec ces
formules.

\begin{omfigure}
\begin{tikzpicture}[font=\small]
\foreach \i/\x in {0/0,1/2.2,2/4.4,3/6.6} {\node[draw=omDef,circle,thick,minimum size=0.85cm,fill=omDef!6] (s\i) at (\x,0) {$\i$};}
\node (dots) at (8.6,0) {$\cdots$};
\foreach \a/\b/\l in {0/1/0,1/2/1,2/3/2} {
  \draw[-{Stealth},thick,omThm] (s\a) to[bend left=35] node[above] {$\lambda_{\l}$} (s\b);
  \pgfmathtruncatemacro{\m}{\l+1}
  \draw[-{Stealth},thick,omDef] (s\b) to[bend left=35] node[below] {$\mu_{\m}$} (s\a);}
\draw[-{Stealth},thick,omThm] (s3) to[bend left=35] node[above] {$\lambda_3$} (dots);
\draw[-{Stealth},thick,omDef] (dots) to[bend left=35] node[below] {$\mu_4$} (s3);
\end{tikzpicture}

\omcaption{Une chaîne de naissance et de mort~: depuis l'état $n$, les seuls mouvements sont vers le haut au taux $\lambda_n$ et vers le bas au taux $\mu_n$. La
meilleure file, comptée en ordres, en est une~: les ordres à cours limité jouent le rôle des naissances, et les ordres au marché et les
annulations celui des morts.}\label{fig:qm:markov-chains-and-queues:bd}
\end{omfigure}

\begin{omfigure}
\begin{tikzpicture}
\begin{axis}[width=11cm,height=4.8cm,
  xlabel={taux d'utilisation $\rho = \lambda/\mu$},ylabel={nombre moyen dans le système},
  tick label style={font=\small},label style={font=\small},
  xmin=0,xmax=0.95,ymin=0,ymax=10.5,grid=major,grid style={gray!25},
  legend columns=3,legend style={font=\footnotesize,at={(0.5,-0.3)},anchor=north,draw=none,fill=none,/tikz/every even column/.append style={column sep=8pt}},
  table/col sep=comma]
\addplot[omDef,very thick] table[x=rho,y=theory]{figdata/methods/08-markov-chains-and-queues/mm1.csv};
\addplot[omThm,only marks,mark=*,mark size=1.8pt] table[x=rho,y=sim]{figdata/methods/08-markov-chains-and-queues/mm1.csv};
\addplot[black,only marks,mark=x,mark size=2.6pt] table[x=rho,y=lamW]{figdata/methods/08-markov-chains-and-queues/mm1.csv};
\legend{$\rho/(1-\rho)$,moyenne temporelle (sim.),$\lambda W$ (sim.)}
\end{axis}
\end{tikzpicture}

\omcaption{La \omterm{def:qm:markov-chains-and-queues:mm1}{file M/M/1}~: nombre moyen dans le système en fonction du taux d'utilisation, à partir de la \omterm{def:qm:stochastic-differential-equations:kolmogorov}{loi stationnaire} et de
200\,000 clients simulés par point, dont le nombre moyen en temps concorde avec $\lambda W$ (loi de Little). Près
de la saturation, la simulation converge lentement. Données~: le tutoriel du chapitre,
avec graine fixée.}\label{fig:qm:markov-chains-and-queues:mm1}
\end{omfigure}

\section{Temps d'atteinte et probabilités d'atteinte}

\begin{definition}[État absorbant, probabilité d'atteinte]\label{def:qm:markov-chains-and-queues:hitting}
Un \emph{état absorbant}\index{etat@état absorbant} est un état que la chaîne ne quitte jamais ($q_i = 0$). Pour un
ensemble $A$ d'états, la \emph{probabilité d'atteinte}\index{probabilité d'atteinte} $h_i = \P_i(\text{la chaîne finit par entrer dans } A)$.
\end{definition}

\begin{method}[Analyse au premier pas]\label{met:qm:markov-chains-and-queues:firststep}
L'\emph{analyse au premier pas}\index{analyse au premier pas} conditionne par le premier saut. Les \omterm{def:qm:markov-chains-and-queues:hitting}{probabilités d'atteinte} résolvent
$h_i = 1$ sur $A$ et $\sum_jq_{ij}h_j = 0$ hors de $A$~; les \omterm{def:qm:brownian-motion:passage}{temps d'atteinte} moyens $m_i$ résolvent $m_i = 0$ sur $A$ et $\sum_jq_{ij}m_j =
-1$ hors de $A$ (les solutions positives ou nulles minimales). Pour une chaîne de naissance et de mort, ces systèmes sont tridiagonaux et
coûtent $O(n)$ à résoudre.
\end{method}

Pour une file de taux constants $\lambda$ et $\mu$, la probabilité d'atteindre $n$ avant 0 depuis $i$ est la
ruine du joueur du chapitre~1, $(1 - (\mu/\lambda)^i)/(1 - (\mu/\lambda)^n)$, et le temps moyen pour vider la file depuis $i$ lorsque $\mu > \lambda$
vaut $i/(\mu - \lambda)$. Les moyennes cachent beaucoup de choses ici~: la loi du temps de vidange est dissymétrique, et une course entre deux
files exige la loi entière. Elle provient de l'uniformisation~: avec $L$ au moins égal à chaque taux total, la chaîne
est une \omterm{def:qm:markov-chains-and-queues:generator}{chaîne des sauts} de matrice de transition $I + \mathcal L/L$, parcourue aux instants d'un \omterm{def:qm:jump-processes:poisson}{processus de Poisson} de taux $L$, de sorte que
$P(t) = \sum_k e^{-Lt}\frac{(Lt)^k}{k!}(I + \mathcal L/L)^k$, somme de termes positifs, stable numériquement à calculer
(tutoriel, premier listing).

\section{Files dimensionnées pour les carnets d'ordres}

On compte chaque meilleure file en ordres de 10 lots. De chaque côté, les ordres à cours limité rejoignent la file au taux $\lambda = 0.9$ par seconde, et
les ordres sortent, par ordres au marché ($0.6$) et annulations ($0.4$), au taux $\nu = 1$~: les files dérivent vers
zéro, et lorsqu'une se vide, le prix bouge. Supposons les deux files indépendantes. Alors le prix milieu monte
au prochain mouvement si le temps de vidange $T_a$ de la file vendeuse survient avant celui de la file acheteuse, et
\[
\P(\text{hausse}) = \P(T_a < T_b) = \int_0^\infty\bigl(1 - F_b(t)\bigr)\,dF_a(t),
\]
où $F_a, F_b$ désignent les lois des temps de vidange. La \Cref{fig:qm:markov-chains-and-queues:moveup} la montre en
fonction de la taille des files~: avec 800 lots à l'achat et 200 à la vente, le prochain mouvement est une hausse avec une probabilité de 95.4~\%. Le
déséquilibre des files prédit le prochain tick, un fait que le livre~10 transforme en signal.

Pour l'ordre du teneur de marché en queue de la file acheteuse, avec $k$ ordres devant lui, le temps d'exécution est le temps que
mettent les $k$ ordres devant à sortir (au taux $\nu$ chacun), augmenté de l'attente d'un ordre au marché (taux $0.6$), indépendamment de
la file vendeuse. La probabilité d'exécution est la course de ce temps contre $T_a$
(voir la \cref{fig:qm:markov-chains-and-queues:fill}). Avec 80 ordres devant et 20 à la vente, elle vaut 73.2~\%, et
73.4~\% dans 20\,000 courses simulées~; avec 20 devant, elle vaudrait 99.3~\%. Une place proche de la tête de la file
vaut de l'argent, ce qui explique que la priorité temporelle soit tant disputée (livre~1, chapitre~19).

\begin{omfigure}
\begin{tikzpicture}
\begin{groupplot}[group style={group size=2 by 1,horizontal sep=1.6cm},
  width=7.6cm,height=4.9cm,tick label style={font=\small},label style={font=\small},
  grid=major,grid style={gray!25},ymin=0,ymax=1.05,table/col sep=comma]
\nextgroupplot[xlabel={file acheteuse (ordres de 10 lots)},ylabel={P(mouvement suivant à la hausse)},xmin=0,xmax=100,
  legend style={font=\footnotesize,at={(0.97,0.03)},anchor=south east,fill=white,draw=none}]
\addplot[omDef,very thick] table[x=bid,y=ask10]{figdata/methods/08-markov-chains-and-queues/moveup.csv};
\addplot[omThm,very thick] table[x=bid,y=ask20]{figdata/methods/08-markov-chains-and-queues/moveup.csv};
\addplot[omProp,very thick] table[x=bid,y=ask40]{figdata/methods/08-markov-chains-and-queues/moveup.csv};
\legend{l'ask 10,l'ask 20,l'ask 40}
\nextgroupplot[xlabel={ordres devant le nôtre},ylabel={P(exécution avant la vidange de l'ask)},ylabel style={font=\footnotesize},xmin=0,xmax=120,
  legend style={font=\footnotesize,at={(0.03,0.03)},anchor=south west,fill=white,draw=none}]
\addplot[omDef,very thick] table[x=ahead,y=ask10]{figdata/methods/08-markov-chains-and-queues/fill.csv};
\addplot[omThm,very thick] table[x=ahead,y=ask20]{figdata/methods/08-markov-chains-and-queues/fill.csv};
\addplot[omProp,very thick] table[x=ahead,y=ask40]{figdata/methods/08-markov-chains-and-queues/fill.csv};
\addplot[black,only marks,mark=o,mark size=2.6pt] coordinates {(80,0.7324)};
\legend{l'ask 10,l'ask 20,l'ask 40}
\end{groupplot}
\end{tikzpicture}

\omcaption{Deux meilleures files traitées comme des chaînes de naissance et de mort indépendantes (arrivées à 0.9, départs à 1 par seconde).
À gauche~: probabilité que le prochain mouvement du prix milieu soit une hausse, en fonction de la file acheteuse, pour trois files vendeuses. À droite~:
probabilité qu'un ordre ayant un nombre donné d'ordres devant lui à l'achat soit exécuté avant que la
file vendeuse ne se vide~; le cercle marque le problème du week-end, 80 devant et 20 à la vente~: 73~\%. Données~: le tutoriel du
chapitre.}\label{fig:qm:markov-chains-and-queues:moveup}\label{fig:qm:markov-chains-and-queues:fill}
\end{omfigure}

\section{Tutoriel~: en queue de file}\label{tut:qm:markov-chains-and-queues:queue}

\begin{tutorial}
\textbf{Objectif.} Calculer la loi du temps de vidange d'une file par uniformisation, la confronter à un temps d'exécution, et
vérifier le résultat par simulation. \textbf{État final~:}
\cref{fig:qm:markov-chains-and-queues:mm1,fig:qm:markov-chains-and-queues:moveup} et la probabilité 73~\%.
\begin{tutsteps}
\item \textbf{Le temps de vidange} d'une file, et la course entre deux temps indépendants.
\omcode{code/firm/queues/firm_queues.py}{104}{128}{Uniformisation d'une file de naissance et de mort absorbée en zéro, et course de deux temps.}\label{lst:qm:markov-chains-and-queues:depletion}
\item \textbf{Le temps d'exécution}~: $k$ départs devant, puis un ordre au marché.
\omcode{code/firm/queues/firm_queues.py}{131}{149}{Temps d'exécution d'un ordre en queue de file.}\label{lst:qm:markov-chains-and-queues:fill}
\item \textbf{Exécuter} \texttt{problem()}, \texttt{spread\_chain()}, \texttt{mm1\_simulation(rho)} et
\texttt{fig\_queues.py}.
\end{tutsteps}
\textbf{Pour aller plus loin.} Rendre les annulations proportionnelles à la taille de la file (taux $\theta$ par ordre) et observer
comment la course se modifie~; faire dépendre les taux d'arrivée du déséquilibre, comme dans les modèles à files réactives du
livre~10.
\end{tutorial}

\section{Construction~: files pour les carnets d'ordres}\label{bld:qm:markov-chains-and-queues:queues}

\begin{build}
\bfield{Rôle} Valoriser la position dans une file, prévoir le prochain mouvement de prix à partir des meilleures files, et
dimensionner les files de messages dans les systèmes de la société.
\bfield{Interface} \texttt{generator(rates)}~; \texttt{stationary(Q)}~; \texttt{jump\_chain(Q)}~;
\texttt{bd\_hitting\_probability(birth, death, n)}~; \texttt{bd\_expected\_exit\_time(birth, death, n)}~;
\texttt{depletion\_cdf(birth, death, start, t)}~; \texttt{fill\_time\_cdf(k, nu, mu, t)}~; \texttt{race(cdf\_a,
cdf\_b, t)}~; \texttt{thomas(a, b, c, d)}.
\bfield{Règles} Taux par seconde, files en ordres~; les taux de naissance et de mort peuvent être des fonctions de l'état pour les
solveurs à premier pas~; uniformisation, et non exponentielles matricielles, pour les lois transitoires~; chaque solveur est confronté
à la simulation dans les tests.
\bfield{Tests d'acceptation} \texttt{code/firm/queues/tests/}~: ruine du joueur et temps de sortie $i(n - i)/(\lambda + \mu)$ dans le
cas symétrique~; \omterm{def:qm:stochastic-differential-equations:kolmogorov}{loi stationnaire} de la M/M/1~; loi de vidange confrontée à la simulation~; course symétrique à files égales.
\bfield{Pour aller plus loin} Taux dépendant de l'état dans la course (une chaîne à deux dimensions résolue par Gauss--Seidel)~; files alimentées
par le flux de Hawkes du chapitre~7.
\end{build}

\begin{omsources}
\item R.~Cont, S.~Stoikov et R.~Talreja, «~A stochastic model for order book dynamics~», \emph{Operations
Research} 58, 2010.
\item R.~Cont et A.~de Larrard, «~Price dynamics in a Markovian limit order market~», \emph{SIAM Journal on
Financial Mathematics} 4, 2013.
\item J.~D.~C.~Little, «~A proof for the queuing formula $L = \lambda W$~», \emph{Operations Research} 9, 1961.
\item D.~G.~Kendall, «~Stochastic processes occurring in the theory of queues and their analysis by the method
of the imbedded Markov chain~», \emph{Annals of Mathematical Statistics} 24, 1953.
\end{omsources}

\section{Exercices}

\begin{exercise}[$\star$]\label{exo:qm:markov-chains-and-queues:1}
Dans la chaîne de la fourchette de l'\cref{ex:qm:markov-chains-and-queues:spread}, quel est le temps moyen passé à chaque fourchette
avant un changement, et vers quel état la chaîne se dirige-t-elle lorsqu'elle quitte 3~ticks~?
\end{exercise}

\begin{exercise}[$\star$]\label{exo:qm:markov-chains-and-queues:2}
Une \omterm{def:qm:markov-chains-and-queues:mm1}{file M/M/1} traite un message par microseconde et en reçoit 0.8 par microseconde. Quel est le nombre moyen
de messages dans le système, et quel est le temps moyen qu'un message y passe~?
\end{exercise}

\begin{exercise}[$\star$]\label{exo:qm:markov-chains-and-queues:3}
Une file de 20 ordres gagne un ordre au taux de 0.9 et en perd un au taux de 1. Quelle est la probabilité qu'elle atteigne 40 avant
de se vider~?
\end{exercise}

\begin{exercise}[$\star\star$]\label{exo:qm:markov-chains-and-queues:4}
La chaîne de la fourchette est-elle réversible~? Vérifier l'\omterm{def:qm:markov-chains-and-queues:balance}{équilibre détaillé} entre 1 et 2~ticks, et le critère de Kolmogorov sur
le cycle $1 \to 2 \to 3 \to 1$.
\end{exercise}

\begin{exercise}[$\star\star$]\label{exo:qm:markov-chains-and-queues:5}
Pour la file de l'\cref{exo:qm:markov-chains-and-queues:3} sans limite supérieure, quels sont le temps moyen et le temps médian
avant qu'elle ne se vide, et la probabilité qu'elle se vide en moins d'une minute~? Pourquoi la moyenne et la médiane diffèrent-elles autant~?
\end{exercise}

\begin{exercise}[$\star\star$]\label{exo:qm:markov-chains-and-queues:6}
Un niveau de prix compte 400 lots en moyenne, et 50 lots par seconde le quittent par exécution ou par annulation. Combien de temps
un lot y reste-t-il en moyenne~?
\end{exercise}

\begin{exercise}[$\star\star\star$]\label{exo:qm:markov-chains-and-queues:7}
\emph{Programmation.} Vérifier la probabilité d'exécution du problème du week-end en simulant la course événement par événement, et
comparer avec le résultat de l'uniformisation.
\end{exercise}

\begin{exercise}[$\star\star\star$]\label{exo:qm:markov-chains-and-queues:8}
\emph{Repérer le défaut.} «~Nos ordres passifs sont exécutés 95~\% du temps lorsque le carnet est à ce point déséquilibré, donc
nous nous plaçons à l'arrière de la file du côté chargé dès que nous le voyons.~»
\end{exercise}

\section{Problème~: l'arrière de la file}

\begin{problem}[{Problème du week-end --- 800 lots devant, 200 de l'autre côté}]\label{pb:qm:markov-chains-and-queues:1}
L'ordre d'un teneur de marché se place à l'arrière du meilleur cours acheteur de 800 lots~; le meilleur cours vendeur en compte 200. Les ordres portent sur 10 lots.
De chaque côté, les ordres à cours limité arrivent au taux de 0.9 par seconde, les ordres au marché au taux de 0.6 par seconde et les annulations au taux de
0.4, donc les ordres quittent la file au taux $\nu = 1$ par seconde. Supposer les files indépendantes.

\medskip\noindent\textbf{Partie I --- Le modèle.}
\begin{enumerate}
\item Combien d'ordres se trouvent devant le nôtre, et du côté vendeur~?
\item Quel est le temps moyen avant que notre ordre ne soit exécuté, en négligeant le côté vendeur~?
\item Quels sont le temps moyen et le temps médian pour que la file vendeuse se vide~?
\item Quelle est la probabilité que la file vendeuse se vide en moins d'une minute~?
\item Quelle est la probabilité qu'elle finisse par se vider~?
\end{enumerate}

\medskip\noindent\textbf{Partie II --- La course.}
\begin{enumerate}[resume]
\item Quelle est la probabilité que notre ordre soit exécuté avant que la file vendeuse ne se vide~?
\item Que donne une simulation de 20\,000 courses~?
\item Et si la file vendeuse comptait 400 ou 800 lots~?
\item Et si seuls 20 ordres étaient devant le nôtre~?
\item Quelle est la probabilité que le prochain mouvement du prix milieu soit à la hausse, avec des files égales~?
\end{enumerate}

\medskip\noindent\textbf{Partie III --- Premiers pas.}
\begin{enumerate}[resume]
\item Quelle est la probabilité que la file vendeuse atteigne 40 ordres avant de se vider~?
\item Quel est le temps moyen jusqu'à ce qu'elle atteigne pour la première fois 0 ou 40~?
\item Comment des annulations proportionnelles à la taille de la file modifieraient-elles la dérive~?
\item Quel effet aurait un flux d'ordres groupés (chapitre~7) sur la queue de la loi du temps de vidange~?
\item Que laisse de côté l'indépendance des deux files~?
\end{enumerate}

\medskip\noindent\textbf{Partie IV --- Jugement.}
\begin{enumerate}[resume]
\item Que vaut une position dans la file, dans ce modèle~?
\item Lorsque l'ordre est effectivement exécuté dans cette configuration, qu'apprend cette exécution au teneur de marché~?
\item Faut-il se placer au cours acheteur, tout compte fait~?
\item Énoncer le \emph{résultat classique}~: la probabilité d'exécution et le temps d'attente moyen.
\item En une phrase~: qu'est-ce qui détermine si un ordre passif placé à l'arrière d'une file est exécuté~?
\end{enumerate}
\end{problem}

\section{Questions d'entretien}

\begin{interviewq}[$\star$ \iqroles{chercheur, développeur}]\label{iq:qm:markov-chains-and-queues:1}
Un serveur traite 10\,000 requêtes par seconde et en reçoit 9\,000. En moyenne, combien de requêtes attendent, et pendant combien de temps~?
\end{interviewq}

\begin{interviewq}[$\star$ \iqroles{chercheur}]\label{iq:qm:markov-chains-and-queues:2}
Qu'est-ce que la \omterm{def:qm:markov-chains-and-queues:generator}{matrice génératrice} d'une \omterm{def:qm:markov-chains-and-queues:ctmc}{chaîne de Markov à temps continu}, et comment en obtenez-vous la \omterm{def:qm:stochastic-differential-equations:kolmogorov}{loi stationnaire}~?
\end{interviewq}

\begin{interviewq}[$\star\star$ \iqroles{trader, chercheur}]\label{iq:qm:markov-chains-and-queues:3}
Le cours acheteur compte 800 lots et le cours vendeur 200. Dans quel sens le prochain tick est-il le plus probable, et pourquoi~?
\end{interviewq}

\begin{interviewq}[$\star\star$ \iqroles{chercheur}]\label{iq:qm:markov-chains-and-queues:4}
Une marche aléatoire équilibrée part de 3. Quelle est la \omterm{def:qm:markov-chains-and-queues:hitting}{probabilité d'atteinte} de 10 avant 0, et quel est le nombre moyen de pas~?
\end{interviewq}

\begin{interviewq}[$\star\star$ \iqroles{trader}]\label{iq:qm:markov-chains-and-queues:5}
Pourquoi les teneurs de marché accordent-ils de la valeur à une place en tête de file, et comment la chiffreriez-vous~?
\end{interviewq}

\begin{interviewq}[$\star\star\star$ \iqroles{développeur, chercheur}]\label{iq:qm:markov-chains-and-queues:6}
Comment calculez-vous la probabilité qu'une chaîne se trouve dans un état donné à l'instant $t$, pour une chaîne comptant des milliers
d'états, de façon numériquement stable~?
\end{interviewq}
